\subsection{Four-model comparison: the metric and ridge intercept change the reading}
The historical comparison on 171 cached Turing simulations reports linear-theory relative RMSE .3427, dispersion-curve CNN .3386, parameter MLP .4487 and curve-feature ridge 1.2396. These exact values remain in \path{results/turing_4model.json}. The .0041 CNN difference is a small point-estimate difference, not an established statistically meaningless effect or a demonstrated improvement: the saved summary does not include paired predictions, intervals or a hypothesis test. Neither a small edge nor an unsuccessful learner proves what information its inputs contain. We withdraw the earlier claims that only linear-theory features matter and that the shift is unlearnable from those views.

The target is $y=(k^2_{sim}-k^2_{min})/(k^2_{max}-k^2_{min})$. The recorded relative RMSE is\[\sqrt{N^{-1}\sum_i ((\hat y_i-y_i)/(|y_i|+10^{-9}))^2},\]\noindent which weights squared target errors by inverse squared target magnitude. This differs from ordinary target RMSE. A target near zero receives much more weight; changing the metric can change the ranking without changing a single prediction. Further, the historical fixed linear predictor is scored globally on all starts, while the trained learners are summarized by the arithmetic mean of twelve fold RMSEs (six folds, two partition seeds). Mean fold root errors and a pooled root error need not agree. Two partitions of the same 171 simulations are not two independent physical replications.

The ridge implementation also standardizes its features within each training fold but fits no intercept. Centered features cannot represent the nonzero mean target using a linear slope alone. That is a comparator specification defect, not evidence that curve summaries add noise rather than signal. A fresh frozen local audit retains identical folds, features, ridge penalty one and targets, and adds an unpenalized intercept using only the training-fold target mean. It includes the original no-intercept version, the fixed linear predictor and a training-mean constant. No CNN or MLP is retrained, and no hyperparameter is selected using these outcomes.

\begin{table}[ht]\centering\small
\caption{Same saved simulation targets and folds: mean over the two partition seeds. Each model has out-of-fold predictions for all 171 simulations per seed. These are descriptive metric comparisons, not independent-replicate uncertainty.}
\begin{tabular}{lrrr}\toprule
Predictor & pooled rel. RMSE & fold-mean rel. RMSE & pooled RMSE\\\midrule
Linear theory & .3427 & .3417 & .1865\\
Training-mean constant & .7381 & .7021 & .1756\\
Ridge, no intercept & 1.2466 & 1.2396 & .4390\\
Ridge, with intercept & .3579 & .3511 & .1372\\\bottomrule
\end{tabular}\end{table}

Adding the intercept reduces mean-fold relative ridge error from 1.2396 to .3511. Under that metric, the repaired ridge still does not beat the fold-matched linear value .3417. Under ordinary pooled target RMSE the same repaired ridge is better (.1372 versus .1865), and even the training-mean constant is better than linear (.1756). Thus a universal ``nothing beats linear theory'' heading is false once the loss is stated differently. This is a useful comparator audit, not a new optimized predictor or an externally validated performance gain. The training loss in the historical neural experiment is ordinary squared error, while its reported score is relative error; this mismatch is another reason not to turn its ranking into an information boundary.

The saved source SHA-256, local plan hash, all per-seed metric records, target range and mean are in \path{results/turing_metric_audit.json}, generated by \path{experiments/turing_metric_audit.py}. Tests reproduce the saved no-intercept result, verify the metric-dependent ranking and regenerate this deterministic audit byte-for-byte. The historical result file and its negative interpretation remain visible as superseded claims, not silently edited measurements. Independent simulations, preselected loss, matched model selection and saved paired neural predictions would be needed for a broader comparison. Nothing here establishes how biological tissues predict or select wavenumber.
